\section{Implementacja komponentu renderującego}
% https://www.open-std.org/jtc1/sc22/wg21/docs/papers/2023/n4950.pdf s. 745
Implementację można podzielić w logiczny sposób na trzy główne etapy: inicializacja okna, inicializajcja WebGPU, oraz pętla główna aplikacji. Etapy inicjalizacji muszą być wykonywane w określonej kolejności, dlatego naturalnie układają się w potok. Jednak inaczej niż w typowym potoku zamiast przekazywać między etapami dane, będziemy przekazywać dostarczony przez standard C++ obiekt \texttt{std::expected<void, std::string>}, symbolizujący stan wykonania poprzedniego etapu: sukces, lub błąd. Następnie etapy zostaną połączone dzięki operacjom na monadach, znanych z programowania funkcyjnego.

\begin{lstlisting}[language=C++, caption={Funkcja inicjalizująca komponent renderujący}, label={lst:initialize}]
void
initialize()
{
  auto result = init_window()
    .and_then([ this ]() { return init_renderer(); })
  
  if (!result.has_value())
  {
    std::println(std::cerr, "Error occured:\n{}", result.error());
    std::exit(1);
  }
}
\end{lstlisting}

\subsection{Okno apliakcji}
Pierwszym etapem jest utworzenie okna aplikacji, w którym będzie wyświetlany renderowany obraz. W tym celu wykorzystano bibliotekę GLFW, która umożliwia tworzenie okien oraz obsługę zdarzeń wejściowych w sposób niezależny od platformy. Proces inicjalizacji okna obejmuje inicializację biblioteki, oraz stworzenie okna. Wskaźnik na utworzone okno jest zapisywany w klasie głównej aplikacji.
\begin{lstlisting}[language=C++, caption={Część kodu odpowiedzialna za inicializację okna}, label={lst:init_window}]
auto
init_window() -> std::expected<void, std::string>
{
  if (glfwInit() == GLFW_FALSE)
  {
    return std::unexpected { "Failed to initialize GLFW" };
  }
  glfwWindowHint(GLFW_CLIENT_API, GLFW_NO_API);
  glfwWindowHint(GLFW_RESIZABLE, GLFW_TRUE);
  p_window =
    glfwCreateWindow(
        initial_width, initial_height, 
        "INZ", nullptr, nullptr);
  if (p_window == nullptr)
  {
    return std::unexpected { "Failed to create window" };
  }
  glfwSetWindowUserPointer(p_window, this);
  
  // ...
}
\end{lstlisting}
Dodatkowo zainicializowane zostały funkcje obsługi zdarzeń, takie jak zmiana rozmiaru okna czy obsługa myszy.
\begin{lstlisting}[language=C++, caption={Dalsza część funkcji inicializującej okno - obsługa pozycji myszy}, label={lst:init_window_cursor}]
auto init_window() -> std::expected<void, std::string>
{
  // ...
  
  glfwSetCursorPosCallback(
    p_window,
    [](GLFWwindow* window, double x_pos, double y_pos)
    {
      if (
        auto* app = 
            reinterpret_cast<application*>(
                glfwGetWindowUserPointer(window)); 
        app != nullptr)
      {
        app->cursor_position_callback(x_pos, y_pos);
      }
    });
  
  // ...
}
\end{lstlisting}

\subsection{WebGPU}
Inicjalizacja WebGPU składa się z trzynastu kroków, które muszą być wykonane w określonej kolejności, aby poprawnie skonfigurować środowisko renderujące. Poniżej przedstawiono poszczególne kroki wraz z opisem ich funkcji.
%init_renderer()
\begin{lstlisting}[language=C++, caption={}, label={}]
\end{lstlisting}

\subsubsection{Adapter}
%https://www.w3.org/TR/webgpu/#adapters
Adapter to reprezentacja fizycznego lub wirtualnego urządzenia GPU dostępnego w systemie. W tym kroku aplikacja wysyła zapytanie do systemu o dostępne adaptery, a następnie wybiera najbardziej odpowiedni na podstawie określonych kryteriów, takich jak wydajność czy obsługiwane funkcje.
%kod init_adapter()
\begin{lstlisting}[language=C++, caption={}, label={}]
\end{lstlisting}

\subsubsection{Device}
%https://www.w3.org/TR/webgpu/#devices
Device to logiczne przedstawienie wybranego adaptera, które umożliwia tworzenie zasobów procesora graficznego oraz wykonywanie operacji renderujących. W tym kroku aplikacja tworzy obiekt Device, który będzie używany do komunikacji z procesorem graifcznym.
%kod init_device()
\begin{lstlisting}[language=C++, caption={}, label={}]
\end{lstlisting}

\subsubsection{Command queue}
%https://eliemichel.github.io/LearnWebGPU/next/getting-started/the-command-queue.html
Skoro Device jest już utworzony, należy teraz stworzyć kolejkę poleceń (Command queue). Dzięki niej kod aplikacji wykonywany na CPU może wysyłać polecenia do GPU w sposób asynchroniczny. W tym kroku aplikacja tworzy obiekt Command queue.
%kod init_queue()
\begin{lstlisting}[language=C++, caption={}, label={}]
\end{lstlisting}

\subsubsection{Shader module}
Podczas tego kroku aplikacja ładuje i kompiluje shadery, które będą używane w trakcie potoku renderującego. Shadery są napisane w języku WGSL i definiują sposób przetwarzania wierzchołków oraz obliczania koloru pikseli.
%kode init_shader_module()
\begin{lstlisting}[language=C++, caption={}, label={}]
\end{lstlisting}

\subsubsection{Depth texture}
Jest to specjalny obiekt do przechowywania informacji o głębi sceny 3D. W tym kroku aplikacja tworzy teksturę głębi, która będzie używana do testów głębi podczas potoku renderującego.
%kod init_depth_texture()
\begin{lstlisting}[language=C++, caption={}, label={}]
\end{lstlisting}

\subsubsection{Bind group layout}
%https://www.w3.org/TR/webgpu/#gpubindgrouplayout
Bind group layout definiuje interfejs między zbiorem zasobów powiązanych w bind group (tworzony na końcu) a ich dostępnością w etapach shadera. Definiowane są tutaj typy zasobów, ich widoczność w poszczególnych shaderach.
%kod init_bind_group_layout()
\begin{lstlisting}[language=C++, caption={}, label={}]
\end{lstlisting}

\subsubsection{Pipeline}
%https://www.w3.org/TR/webgpu/#gpurenderpipeline
Głównym elementem procesu renderującego jest potok renderujący (Render Pipeline). W tym kroku aplikacja tworzy obiekt Pipeline, który definiuje wszystkie etapy przetwarzania danych. Korzystając z oficjalnego standardu WebGPU, można wyszczególnić następujące etapy etapy potoku renderującego:
\begin{enumerate}
    \item Vertex fetch                  % Pobieranie wierzchołków
    \item Vertex shader                 % =
    \item Primitive assembly            % Prymitywny Assembly
    \item Rasterization                 % Rasteryzacja
    \item Fragment shader               % =  
    \item Stencil test and operation    % ???
    \item Depth test and write          % Test głębi i zapis??? 
    \item Output merging                % Scalanie przetworzonych danych
\end{enumerate}
\subsubsection{Sampler}
\subsubsection{Textures}
\subsubsection{Geometry}
\subsubsection{Uniforms}
\subsubsection{Lighting uniforms}
\subsubsection{Bind group}

\subsection{Pętla główna}
